% HUONG DAN SU DUNG SACH
\chapter*{Hướng Dẫn Sử Dụng Sách}
\addcontentsline{toc}{chapter}{Hướng Dẫn Sử Dụng Sách}
\markboth{Hướng Dẫn Sử Dụng Sách}{Hướng Dẫn Sử Dụng Sách}

\begin{truyen}{Quyển Sách Này Dành Cho Ai?}{Who Is This Book For?}
\chuhoa{Q}{uyển sách này phù hợp cho ba nhóm người đọc.}
\textit{(This volume is suitable for three groups of readers.)}

\textbf{Học sinh:} để thấy nhiều câu ``em chỉ nói cho vui'' thật ra chạm vào chuyện rất nghiêm túc của việc học.
\textit{(\textbf{Students:} to see that many funny excuses actually touch serious matters of learning.)}

\textbf{Giáo viên:} để lấy chất liệu mở đầu tiết học, sinh hoạt chủ nhiệm, hoặc trò chuyện cuối giờ.
\textit{(\textbf{Teachers:} to use as material for lesson openings, class meetings, or end-of-period conversations.)}

\textbf{Phụ huynh:} để hiểu hơn tâm lý con em tuổi đi học và có thêm cách nói chuyện nhẹ mà sâu.
\textit{(\textbf{Parents:} to better understand school-age children and find gentler but deeper ways to talk with them.)}
\end{truyen}

\vspace{6pt}

\begin{tcolorbox}[
  colback=sunclay!8,
  colframe=sunclay!70!black,
  coltitle=black,
  fonttitle=\bfseries,
  title={Giải Thích Các Ô Trong Sách \quad\textit{\small(Book Guide)}},
  arc=4pt, boxrule=1pt, breakable, enhanced
]

\smallskip
\textbf{Ô Xanh Đậm --- Màn Đối Đáp Chính (\textit{Main Dialogue})}\\
Phần này là câu hỏi lý sự của học sinh và cách người thầy đáp lại. Hãy đọc chậm để thấy sự chuyển từ chống chế sang suy nghĩ.\\
\textit{(This is the student excuse and the teacher reply. Read slowly to notice the shift from resistance to reflection.)}

\medskip
\textbf{Ô Đen Xám --- Bài Học Cuối Giờ (\textit{Takeaway})}\\
Một hoặc hai câu ngắn chốt lại điều cần giữ. Đây là phần nên đọc lại nhiều lần.\\
\textit{(One or two short sentences that keep the main lesson. This is the part worth rereading.)}

\medskip
\textbf{Ô Xanh Nhạt --- Easy English}\\
Hai câu tiếng Anh ngắn để nhớ nhanh ý chính và luyện thêm từ vựng.\\
\textit{(Two short English lines to remember the core idea and practise vocabulary.)}

\medskip
\textbf{Ô Vàng --- Gợi Ý Tự Nghĩ (\textit{Self-Reflection})}\\
Hai câu hỏi rất ngắn để người đọc soi lại chính mình. Có thể trả lời bằng miệng, viết ra, hoặc đem ra thảo luận trong lớp.\\
\textit{(Two brief reflection questions for personal thinking, writing, or class discussion.)}

\smallskip
\end{tcolorbox}

\vspace{6pt}

\begin{baihoc}
Gợi ý cách đọc hiệu quả nhất: mỗi ngày 1 đến 2 đoạn. Đọc xong thì dừng lại, đừng lướt qua quá nhanh vì tiếng cười thường đi trước phần đáng nghĩ.
\textit{(Best reading method: one or two pieces a day. Pause after reading, because laughter often comes before the real point.)}
\end{baihoc}

\begin{ghinhoanh}
\textbf{Reading tips:}

Read slowly. Laugh first, then think.

One short dialogue can hold one big lesson.
\end{ghinhoanh}

\ngancach